\section{Method}
\begin{frame}{Assumptions and Approximations}
    \begin{block}<+->{Assumptions}
        \begin{itemize}
            \item<+-> No planet migration
            \item<+-> Geochemical gradient extrapolation
            \item<+-> $T>T_{\mathrm{cond}}(\mathrm{CaS},\mathrm{MgS})$ %(fit Ca:Al, Ca:Ti, Al:Ti observations)
        \end{itemize}
    \end{block}
    \begin{block}<+->{Shell v Core}
        \centering
            \begin{tabular}{c|c|c}
            Element & Shell & Core \\
            \hline
            \onslide<+->{Mg & $\checkmark$ & - \\
            Al & $\checkmark$ & - \\
            Ca & $\checkmark$ & - \\}
            \onslide<+->{Ni & $40\%$ & $60\%$ \\
            S & $40\%$ & $60\%$ \\
            Si & $65\%$ & $35\%$}
            \end{tabular}
    \end{block}
\end{frame}
\begin{frame}{Model Inputs}
    \begin{block}{Data sources}
        \centering
        \begin{tabular}{c|c}
            Source & Data \\
            \hline
            \hline
            \onslide<+->{Mercury surface & K:Th \\
                            & Mg\#} \\
            \hline
            \onslide<+->{E-chondrites & Ca:Al \\
                        & Al:Ti \\
                         & Fe:Ni} \\
            \hline
            \onslide<+->{Earth & Mg:Al} \\
            \hline
            \onslide<+->{Extrapolation & Fe:Si \\
                          & Mg:Si \\
                          & Al:Si\\
                          & S} \\
        \end{tabular}
    \end{block}
        % \begin{block}{Data sources}
        %     \begin{itemize}
        %         \item \textbf{Mercury surface} K:Th, Mg#
        %         \item \textbf{E-chondrites} Ca:Al, Fe:Ni
        %         \item \textbf{Earth} Mg:Al 
        %         \item \textbf{Extrapolation} Fe:Si, Mg:Si, Al:Si
        %     \end{itemize}
        % \end{block}
\end{frame}

% \begin{frame}{Model inputs}
%     % \begin{block}{Atomic Abundances}
%         % \begin{itemize}
%         %     \item High Fe
%         %     \item Low O
%         %     \item High Fe:O
%         % \end{itemize}
%     % \end{block}
%     \input{FigTex/McDonough21_Fig1}
% \end{frame}

% \begin{frame}{Model inputs}
%     % \begin{block}{Major Cation Ratios}
%     % \begin{itemize}
%     %     \item Extrapolation
%     %     \item Uncompressed density
%     %     \item Si, (Fe:Ni:Si)$_{\mathrm{met}}$, FeS
%     % \end{itemize}
%     % \end{block}
%     \input{FigTex/McDonough21_Fig2}
% \end{frame}

\begin{frame}{Model Inputs}
    % \begin{block}{Lithophile Siderophile Trends}
    %     \begin{itemize}
    %         \item Assume 
    %     \end{itemize}
    % Constrain core-S via extrapolation

    % \end{block}
    \input{FigTex/McDonough21_Fig3}
\end{frame}

\begin{frame}{Model Inputs}
    \input{FigTex/McDonough21_Fig4}
\end{frame}

\begin{frame}{Magnetic Field}
% \begin{block}<+->{\textbf{\input{Equations/FieldStrength}} (\cite{Wardle07})}
%     \begin{itemize}
%         \item \textbf{$\sqrt{|B_z B_\phi|_s}$} Root mean square magnetic field strength
%         \item \textbf{$r_{\mathrm{au}}$} Heliocentric distance
%         \item \textbf{$\dot{M}$} Accretion rate
%         \item \textbf{$v_K$} Local Keplerian speed
%         \item \textbf{$\mathrm{G}$} Gravitational constant
%     \end{itemize}
% \end{block}
\begin{block}{Magnetic field strength {\tiny (\cite{Wardle07})}}
    \begin{itemize}
        \item Needed: $2\ \mathrm{mT}$
        \item Estimate: $\sim 0.3-3\ \mathrm{mT}$
    \end{itemize}
\end{block}
\end{frame}